% Generated draft: requires mathematical and source review.
\documentclass[11pt]{article}
\usepackage[margin=25mm]{geometry}
\usepackage{fontspec}
\setmainfont{Latin Modern Roman}
\usepackage{amsmath,amssymb,mathtools,unicode-math}
\setmathfont{Latin Modern Math}
\usepackage{enumitem,needspace,xurl,hyperref}
\setlength{\parindent}{0pt}
\setlength{\parskip}{5pt}
\setlength{\emergencystretch}{3em}
\newcommand{\N}{\mathbb{N}}
\newcommand{\Z}{\mathbb{Z}}
\newcommand{\Q}{\mathbb{Q}}
\newcommand{\R}{\mathbb{R}}
\newcommand{\C}{\mathbb{C}}
\newcommand{\F}{\mathbb{F}}
\newcommand{\A}{\mathbb{A}}
\newcommand{\PP}{\mathbb{P}}
\newcommand{\E}{\mathbb{E}}
\newcommand{\eps}{\varepsilon}
\newcommand{\dd}{\,\mathrm{d}}
\begin{document}
\section*{UnsolvedMath ID: 2100706}
\textbf{Source problem code:} catalogue entry
\subsection{Definitions and mathematical statement}
A formal power series in the semiclassical parameter is an expression $\mathcal B(\hbar)=\sum_{j\geq0}b_j\hbar^j$; no convergence is implied. In a Bohr--Sommerfeld quantization problem the coefficients can depend on spectral/base variables and on choices of local actions and normalization. The source asks about a particular \emph{singular} Bohr--Sommerfeld series but supplies neither its formula nor the class of singularities and quantizations to which it belongs.

\textbf{Question attributed to S. V\~u Ng\d{o}c.} Is the singular Bohr--Sommerfeld formal power series in $\hbar$ a spectral invariant?

To state the invariant property abstractly, let $P$ denote an admissible quantum integrable system, let $\operatorname{Spec}(P)$ denote the intended spectral data, and let $\mathcal B_P$ denote its specified singular Bohr--Sommerfeld series. The question asks whether
\[
 \operatorname{Spec}(P)=\operatorname{Spec}(Q)
 \quad\Longrightarrow\quad
 \mathcal B_P=\mathcal B_Q,
\]
with equalities replaced by the intended asymptotic or normalization equivalences if those are part of the setting. It asks whether the spectrum determines the series, not whether equality of series determines the whole quantum system.

The record does not specify whether the data are a single spectrum or a joint spectrum of commuting operators, exact eigenvalues or semiclassical data modulo errors such as $O(\hbar^\infty)$, local spectral windows or global spectra, or spectra with multiplicity. It also does not give a coefficient space, action normalization, or equivalence relation for the series. These missing choices are essential and are not silently fixed here.

The supplied research summary reports established connections between singular Bohr--Sommerfeld series and spectral invariants in semitoric/quantum literature, but says the full spectral-invariance statement was not verified as settled. It supplies no precise theorem or reference resolving the question.
\subsection{Short English statement}
Does the relevant quantum spectrum determine the singular Bohr--Sommerfeld formal series, coefficient by coefficient or up to its prescribed normalization? The source omits the singularity class, quantization framework, spectral-data convention, and series normalization, so the exact invariance assertion needs source context.
\subsection{Sources}
\begin{itemize}
\item Local UnsolvedMath registry, permanent ID 2100706, source code catalogue entry, ``Open Problems in Integrable Systems --- Integrability and Quantisation''; attribution supplied with damaged accent markup, rendered here as S. V\~u Ng\d{o}c. Supplied status partially solved, classification PARTIAL-PROGRESS, updated 2026-08-06.
\item \emph{Open Problems, Questions, and Challenges in Finite-Dimensional Integrable Systems} (2018), Problem 7.6, \url{https://arxiv.org/abs/1804.03737}. The record reports source-TeX extraction and access on 2026-07-29. It says the source presents an open item and current status requires release review. Rights: NEEDS\_REVIEW; public arXiv source TeX inspected according to the record.
\item The supplied research summary mentions semitoric/quantum literature and established connections but supplies no identifiable paper or theorem for the partial progress or the full assertion.
\end{itemize}
\end{document}
